\section{Limitations and Differences from the Original Paper}

This implementation adapts the paper's prediction mechanisms to
gem5 MinorCPU rather than reproducing its complete architecture
and evaluation environment~\cite{ravenel2025architecting}. The
differences in processor organization, prediction timing, and
evaluation methodology affect the resulting behavior and
performance.

\subsection{Processor Model and Evaluation Scope}

The implementation uses gem5 MinorCPU and targeted RISC-V
microbenchmarks, whereas the paper uses a modified CVA6-based
processor and SPEC CPU2017 SimPoint evaluation. The paper's
processor is four-wide, with a 32-entry scoreboard, and can issue two
loads per cycle~\cite{ravenel2025architecting}. The MinorCPU
configuration used here is also four-wide but issues at most one memory
instruction per cycle, and its scoreboard tracks every architectural
register rather than a fixed number of entries.

Consequently, the evaluations differ in processor organization,
pipeline constraints, cache configurations, and workloads. The
reported cycle reductions characterize the tested simulator
configurations and deterministic microbenchmarks, not general
application-level speedups, and should not be compared directly
with the paper's performance results.

The value-cache features also support only a single CPU, a single
hardware thread, and syscall-emulation mode, and virtual synonyms,
uncacheable accesses, and compressed loads were not evaluated.

\subsection{Differences from the Original Design}

The LVP predictor table and its confidence policy follow the paper: a
direct-mapped, untagged, PC-indexed table with a two-bit saturating
counter~\cite{ravenel2025architecting}.

The differences lie primarily in prediction use and recovery. The
implementation consumes an LVP prediction only as the base operand of
an eligible dependent integer load. The consumer must be the oldest
unsent memory operation, while its producer must be the oldest
in-flight instruction, carry a prediction, have already entered the
LSQ, and still be awaiting its response. The prediction is validated against the producer's
actual value when its response arrives. If an incorrect prediction was
used, the implementation redirects execution and refetches younger
instructions rather than reproducing the paper's scoreboard-based
replay mechanism.

The ASTRA predictor retains full 64-bit base addresses. The paper's
Last Address Compressed Table (LACT) is a 32-entry LRU address
compressor: a compact 5-bit identifier represents the upper 44 bits,
while the lower 20 bits are retained in the predictor entry. The
paper reports an 85\% reduction in predictor storage from this
representation~\cite{ravenel2025architecting}. The implementation
does not include this compression mechanism. Similarly, the paper's value-cache tags reuse the region
identifier with 8 additional bits, whereas the implementation uses
full virtual-page-number tags. The paper reports a 75\% tag-storage
reduction from its compressed representation
~\cite{ravenel2025architecting}. Although the implemented value cache has the same 4-KiB capacity and 16-byte entry
geometry, it uses per-byte validity information, does not implement a
16-byte line-fill mechanism, and invalidates or falls back for certain
special and cross-line accesses.

V-ASTRA candidacy is decided in Execute when the load initiates its
memory access. At this point, the predicted base is compared with the
actual base and older in-flight stores are scanned for potential
conflicts. This candidate decision does not immediately authorize a
bypass. The final bypass is accepted in the LSQ only after
address translation and the store-buffer check. The implementation
then checks the effective address, repeats the older-store scan using
the actual access size, verifies the saved value-cache snapshot, and
performs a fresh value-cache lookup. If any check fails, the load falls
back to the normal memory path. In contrast, the paper performs
address validation at issue, allowing eligible dependents to issue in
the same cycle as the load.

For store safety, the implementation uses explicit scans of older
in-flight stores and atomics, together with the existing store-buffer
checks. The scan is performed both when forming a V-ASTRA candidate
and when the request is considered for bypass. This differs from the
paper's counting Bloom filter (CBF) and in-flight invalidation FIFO,
which track potential conflicts with outstanding stores
~\cite{ravenel2025architecting}. The exact
scan avoids the CBF's set-granularity false positives, although its
cost grows with the number of in-flight stores. The FIFO is not
required by the current MinorCPU path because older memory references
are issued in order and their store addresses are known by the
candidate decision point.

The implementation also omits the paper's L0 minicache. The paper
describes a 32-entry direct-mapped VIPT load/store-allocate structure
with separate speculative and committed data arrays, per-byte
validity, and line locking from execution through commit. It provides
a fast path for loads affected by in-flight stores and supports
speculative store-to-load forwarding~\cite{ravenel2025architecting}. In the
current implementation, loads rejected because of older-store conflicts fall back to the
normal memory path. MinorCPU's existing store buffer can forward fully
covered loads, while partial coverage must wait. The minicache would
therefore add functionality beyond the current value cache and bypass
path.

The dependency mechanisms also differ from the paper's early-execution
model. Design~A releases the scoreboard dependency associated with a
bypassed load, allowing a dependent instruction to become eligible
earlier under the stall-on-use timing configuration. It does not
forward a predicted value: the dependent still reads the committed
register value. Under the default timing configuration, this release
mechanism is inactive. Design~B forwards a validated V-ASTRA value as
the base operand of an eligible dependent integer load. It does not
provide general forwarding to arbitrary ALU, branch, store, or
floating-point instructions. Non-memory dependent instructions
continue to execute at in-order commit, after their producer commits.

Both designs recover
only when a wrong prediction was actually used. A wrong prediction that
was never consumed requires no recovery in either. The difference lies
in the recovery mechanism: the paper replays younger instructions
through its scoreboard, whereas this implementation squashes younger
instructions and redirects execution to refetch them.

\subsection{Future Work}

Future work covers project-specific extensions, components from the
original design, and broader evaluation.

\paragraph{Further development of project-specific extensions.}
The random-base result motivates an adaptive admission policy for
provisional LVP-derived bases, potentially using LVP confidence,
recent validation outcomes, or agreement with ASTRA predictions.
It should be evaluated on random base, stable flip, V-cache flip, mixed,
and forced-misprediction workloads to assess gains and regressions.
Further evaluation should isolate Design~B by comparing V-ASTRA
with forwarding enabled and disabled, and measure Design~A under
stall-on-use timing. Design~B could also be extended to additional
load operand positions and store addresses. Forwarding to ALU and
branch instructions would require more substantial changes to
MinorCPU's execution model.

\paragraph{Architectural extensions toward the original design.}

The most substantial omitted component is the speculative L0
minicache. It could provide a fast path for store-dependent loads
currently rejected by V-ASTRA, using speculative and committed data
arrays with store-to-load forwarding. Its implementation would
require LSQ integration, speculative-store handling, line locking,
and recovery on stream changes. Its benefit should be evaluated
using store-heavy workloads and in combination with META.

Another extension is LACT-based address compression. The paper's
32-entry LRU compressor maps 5-bit region identifiers to upper
address bits while retaining 20 low bits in LAT entries. The
value-cache tag reuses the identifier with 8 additional bits. The
paper reports 85\% lower predictor storage and 75\% lower tag
storage. The current implementation instead stores full 64-bit
bases and full virtual-page tags. Implementing LACT would require
identifier allocation and recycling, compressed LAT and value-cache
tags, cache invalidation on identifier recycling, and rejection of
predictions whose displacement carries out of the compressed
region. This is primarily a storage-efficiency and architectural
fidelity extension.

Other possible changes include replacing the exact older-store scan
with the paper's counting Bloom filter (CBF) and invalidation FIFO,
and replacing squash-and-refetch recovery with scoreboard replay.
These are lower priority: the current scan is exact, while replay
would require substantial changes to MinorCPU's reissue and
scoreboard-restoration mechanisms. Further directions discussed
in the paper include VTAGE, process-ID-tagged value-cache entries,
multiprocessor support, and security analysis
\cite{ravenel2025architecting}.

\paragraph{Broader performance evaluation.}

The current evaluation uses deterministic microbenchmarks and a
single main processor and memory configuration. Future work should
include application-level and embedded benchmarks, broader cache
and memory-latency configurations, and different processor timing
parameters and widths where supported. Supporting compressed
RISC-V loads would also improve coverage of binaries built for
\texttt{rv64gc}. Multicore and full-system support would require
additional coherence and external-writer mechanisms, making them
larger extensions beyond the current scope.
